\documentclass[10pt,a4paper]{amsart}
\usepackage[margin=19mm]{geometry}
\usepackage{amsmath,amssymb,mathtools}
\usepackage[scr=rsfs,cal=euler]{mathalfa}
\usepackage{xfrac}
\usepackage[unicode,hidelinks,bookmarks=false]{hyperref}
\newcommand{\ie}{i.e.\ }
\newcommand{\cf}{cf.\ }
\newcommand{\resp}{respectively\ }
\newcommand{\Subsection}{\subsection}
\providecommand{\nobreakdash}{\nobreak\hbox{-}}
\setlength{\emergencystretch}{2em}
\multlinegap0pt
\usepackage{amssymb}
\newtheorem{proposition}{Proposition}
\newtheorem{assumption}{Assumption}
\newtheorem{theorem}{Theorem}
\newtheorem{lemma}{Lemma}
\newcommand{\Be}{\widetilde{B}}       % augmented mobility
\newcommand{\D}{\mathcal{D}}
\newcommand{\Hc}{H_c}                 % composition subspace
\newcommand{\Pc}{P_c}                 % composition projector
\newcommand{\Pm}{P_m}                 % mass projector
\newcommand{\R}{\mathbb{R}}
\newcommand{\bone}{\mathbf{1}}
\title{Entropy Geometry and Augmented Mobility for Reactive Maxwell--Stefan Membrane Transport with Finite Occupancy}
\author{E. Yu. Shchetinin, A. A. Shevchuk, S. I. Salpagarov}
\date{Source: arXiv:2606.22624v1 (CC BY 4.0)}
\begin{document}
\maketitle

\noindent\emph{Source and licence.} Entropy Geometry and Augmented Mobility for Reactive Maxwell--Stefan Membrane Transport with Finite Occupancy. Version arXiv:2606.22624v1 (CC BY 4.0), distributed by arXiv under the Creative Commons Attribution 4.0 International licence, http://creativecommons.org/licenses/by/4.0/, as recorded in the arXiv licence metadata for this exact version and shown on the abstract page https://arxiv.org/abs/2606.22624v1. Full licence notice: \texttt{licenses/arxiv-2606.22624-CC-BY-4.0.txt}.\par\smallskip

\noindent\textit{Editorial scope.} Counted unit: the article's theorem thm:coer (source0.tex lines 456--469). The article's complete original proof of it is delivered with it (source0.tex lines 553--570, one \texttt{proof} environment of 18 lines, which the article places away from the statement). No mathematics is written, repaired or re-derived: the statement and the proof are the article's own lines, in the article's own notation, and no retained line is substituted or re-flowed.  Retained uncounted as the source-local context the proof needs: lines 220--223; lines 230--234; lines 337--350; lines 345--347; lines 373--391; lines 399--429; lines 539--546.  Nothing was omitted from any retained span, so the package carries no omission record.  No retained line contains a citation, so the wrapper carries no bibliography and no bibliography entry.  No retained line contains a figure environment, an image-inclusion command or drawing code, so no image is delivered and the package declares no asset dependency.  Editorial formatting only, in addition: an AMS \texttt{amsart} preamble that restates the article's own declarations and macros used by the retained lines, namely the environments \texttt{proposition}, \texttt{assumption}, \texttt{theorem}, \texttt{lemma} on separate counters, together with the standard packages those lines need.  The retained environments print in this document as Proposition 1, Assumption 1, Theorem 1, Lemma 1 in order of appearance, whereas the article prints the same items with the numbering of its own full document.  The complete native source of the article as distributed in the arXiv e-print for this exact version is bundled unchanged as \texttt{sources/203365/source0.tex}.
\begin{equation}\label{eq:h}
  h(u) = \sum_{i=1}^n \bigl(u_i\log u_i - u_i\bigr)
        + \bigl((1-\varrho(u))\log(1-\varrho(u)) - (1-\varrho(u))\bigr).
\end{equation}

\begin{equation}\label{eq:w}
  w = \nabla h(u), \qquad
  w_i = \partial_{u_i}h(u) = \log u_i - \log(1-\varrho(u)),
  \quad i=1,\dots,n.
\end{equation}

\begin{proposition}\label{prop:augm}
For every $u\in\D$, the rank-one augmentation acts only in the mass
channel. More precisely,
\begin{equation*}
  (\bone\otimes\bone)\Pc = 0, \qquad
  (\bone\otimes\bone)\Pm = n\Pm,
\end{equation*}
and therefore
\begin{equation}\label{eq:Be}
  \Be(u) = B_c(u)\Pc + \bigl(b_m(u)+n\gamma_0\bigr)\Pm.
\end{equation}
In particular, $\Pc\Be(u)z = \Pc B(u)z$ for every $z\in\R^n$, so the
augmentation does not modify the composition block.
\end{proposition}

\begin{equation}
  \Be(u) = B_c(u)\Pc + \bigl(b_m(u)+n\gamma_0\bigr)\Pm.
\end{equation}

Let $\Omega\subset\R^d$ be a bounded domain with Lipschitz boundary,
let $T>0$, and let $u:\Omega\times(0,T)\to\D$ be the unknown state.
The finite-occupancy reactive membrane system studied in this paper is
\begin{subequations}\label{eq:sys}
\begin{alignat}{2}
  \partial_t u - \operatorname{div}\bigl(\Be(u)\nabla w(u)\bigr)
    &= Qu - S(u)
    &&\quad\text{in }\Omega\times(0,T),
    \label{eq:pde}\\
  \Be(u)\nabla w(u)\cdot\nu
    &= 0
    &&\quad\text{on }\partial\Omega\times(0,T),
    \label{eq:bc}\\
  u(\cdot,0) &= u^0
    &&\quad\text{in }\Omega,
    \end{alignat}
\end{subequations}
where $\nu$ denotes the outer unit normal to $\partial\Omega$ and
$w=\nabla h(u)$ is the entropy variable associated with~\eqref{eq:h}--\eqref{eq:w}.

\begin{assumption}\label{ass}
The following hypotheses are in force throughout the paper.
\begin{description}
  \item[(A1)] The initial datum satisfies
    $u^0\in L^\infty(\Omega;\D)$ and
    $\operatorname{ess\,inf}_{x\in\Omega}(1-\varrho(u^0(x)))>0$.
  \item[(A2)] The map $B_c:\D\to\R^{n\times n}$ is continuous,
    symmetric, and nonnegative on $\Hc$. Moreover, there exists
    $\beta_c>0$ such that
    \[
      z^\top B_c(u)z \ge \beta_c|z|^2
      \quad\text{for all }u\in\D,\; z\in\Hc.
    \]
  \item[(A3)] The mass-channel coefficient $b_m:\D\to[0,\infty)$ is
    continuous. No strictly positive lower bound on $b_m$ is assumed.
  \item[(A4)] The reaction matrix $Q\in\R^{n\times n}$ satisfies
    $q_{ij}\ge 0$ for $i\ne j$ and $\bone^\top Q=0$. Hence $Q$
    generates internal conversion without net creation of total mobile
    mass.
  \item[(A5)] The sink term $S:\D\to[0,\infty)^n$ is continuous on
    $\D$ and locally Lipschitz in $\D$. In addition,
    $S_i(u)=0$ whenever $u_i=0$, so that the positive cone is
    invariant under the sink dynamics.
  \item[(A6)] The mobility in~\eqref{eq:pde} is the augmented
    projected mobility
    \[
      \Be(u) = B_c(u)\Pc + b_m(u)\Pm + \gamma_0\,\bone\otimes\bone,
      \qquad \gamma_0>0.
    \]
\end{description}
\end{assumption}

\begin{lemma}\label{lem:decomp}
For every $z\in\R^n$,
\[
  z = \Pc z + \Pm z,\quad
  |z|^2 = |\Pc z|^2 + |\Pm z|^2,\quad
  \Pm z = \tfrac{1}{n}(\bone\cdot z)\,\bone.
\]
\end{lemma}

\begin{theorem}[Channel-wise coercivity of the augmented mobility]
\label{thm:coer}
Let Assumption~\ref{ass} hold. Then there exists a constant $c_*>0$
such that for all $u\in\D$ and all $z\in\R^n$,
\begin{equation}\label{eq:coer}
  z^\top\Be(u)z \ge \beta_c|\Pc z|^2 + n\gamma_0|\Pm z|^2 \ge c_*|z|^2.
\end{equation}
Consequently,
\begin{equation}\label{eq:coer_int}
  \int_\Omega \nabla w(u):\Be(u)\nabla w(u)\,dx
  \ge \beta_c\int_\Omega|\Pc\nabla w(u)|^2\,dx
    + n\gamma_0\int_\Omega|\Pm\nabla w(u)|^2\,dx.
\end{equation}
\end{theorem}

\begin{proof}[Proof of Theorem~\ref{thm:coer}]
Using Proposition~\ref{prop:augm} and~\eqref{eq:Be}, we write
$\Be(u) = B_c(u)\Pc + (b_m(u)+n\gamma_0)\Pm$.
Hence, for any $z\in\R^n$, Lemma~\ref{lem:decomp} gives
\begin{align*}
  z^\top\Be(u)z
  &= (\Pc z)^\top B_c(u)(\Pc z) + (b_m(u)+n\gamma_0)|\Pm z|^2 \\
  &\ge \beta_c|\Pc z|^2 + n\gamma_0|\Pm z|^2 \\
  &\ge \min(\beta_c,n\gamma_0)\bigl(|\Pc z|^2+|\Pm z|^2\bigr)
   = c_*|z|^2,
\end{align*}
which establishes~\eqref{eq:coer} with $c_* := \min(\beta_c,n\gamma_0)>0$; here we used
Assumption~\ref{ass}(A2) for the composition term and $b_m\ge 0$
from~(A3) for the mass term.
The integrated estimate~\eqref{eq:coer_int} follows by applying the
pointwise bound to $z = \partial_{x_\ell}w(u)$ and summing over
$\ell=1,\dots,d$.
\end{proof}

\end{document}
